\documentclass[10pt,a4paper]{amsart}
\usepackage{amsmath,amssymb,amsthm,graphicx}
\usepackage{verbatim}
\newenvironment{reference}{\comment}{\endcomment}
%\newenvironment{reference}{}{}
\newenvironment{slogan}{\comment}{\endcomment}
\newenvironment{history}{\comment}{\endcomment}
% For commutative diagrams we use Xy-pic
\usepackage[all]{xy}

% We use 2cell for 2-commutative diagrams.
\xyoption{2cell}
% Theorem environments.
%
\theoremstyle{plain}
\newtheorem{theorem}[subsection]{Theorem}
\newtheorem{proposition}[subsection]{Proposition}
\newtheorem{lemma}[subsection]{Lemma}

\theoremstyle{definition}
\newtheorem{definition}[subsection]{Definition}
\newtheorem{example}[subsection]{Example}
\newtheorem{exercise}[subsection]{Exercise}
\newtheorem{situation}[subsection]{Situation}

\theoremstyle{remark}
\newtheorem{remark}[subsection]{Remark}
\newtheorem{remarks}[subsection]{Remarks}

\numberwithin{equation}{subsection}

% Macros
%
\def\lim{\mathop{\mathrm{lim}}\nolimits}
\def\colim{\mathop{\mathrm{colim}}\nolimits}
\def\Spec{\mathop{\mathrm{Spec}}}
\def\Hom{\mathop{\mathrm{Hom}}\nolimits}
\def\Ext{\mathop{\mathrm{Ext}}\nolimits}
\def\SheafHom{\mathop{\mathcal{H}\!\mathit{om}}\nolimits}
\def\SheafExt{\mathop{\mathcal{E}\!\mathit{xt}}\nolimits}
\def\Sch{\mathit{Sch}}
\def\Mor{\mathop{\mathrm{Mor}}\nolimits}
\def\Ob{\mathop{\mathrm{Ob}}\nolimits}
\def\Sh{\mathop{\mathit{Sh}}\nolimits}
\def\NL{\mathop{N\!L}\nolimits}
\def\CH{\mathop{\mathrm{CH}}\nolimits}
\def\proetale{{pro\text{-}\acute{e}tale}}
\def\etale{{\acute{e}tale}}
\def\QCoh{\mathit{QCoh}}
\def\Ker{\mathop{\mathrm{Ker}}}
\def\Im{\mathop{\mathrm{Im}}}
\def\Coker{\mathop{\mathrm{Coker}}}
\def\Coim{\mathop{\mathrm{Coim}}}

% Boxtimes
%
\DeclareMathSymbol{\boxtimes}{\mathbin}{AMSa}{"02}

%
% Macros for moduli stacks/spaces
%
\def\QCohstack{\mathcal{QC}\!\mathit{oh}}
\def\Cohstack{\mathcal{C}\!\mathit{oh}}
\def\Spacesstack{\mathcal{S}\!\mathit{paces}}
\def\Quotfunctor{\mathrm{Quot}}
\def\Hilbfunctor{\mathrm{Hilb}}
\def\Curvesstack{\mathcal{C}\!\mathit{urves}}
\def\Polarizedstack{\mathcal{P}\!\mathit{olarized}}
\def\Complexesstack{\mathcal{C}\!\mathit{omplexes}}
% \Pic is the operator that assigns to X its picard group, usage \Pic(X)
% \Picardstack_{X/B} denotes the Picard stack of X over B
% \Picardfunctor_{X/B} denotes the Picard functor of X over B
\def\Pic{\mathop{\mathrm{Pic}}\nolimits}
\def\Picardstack{\mathcal{P}\!\mathit{ic}}
\def\Picardfunctor{\mathrm{Pic}}
\def\Deformationcategory{\mathcal{D}\!\mathit{ef}}
\usepackage{cleveref}
\crefname{originallemma}{original lemma}{original lemmas}
\Crefname{originallemma}{Original lemma}{Original lemmas}
\makeatletter
\newlabel{lemma-base-change-quasi-compact}{{lemma-base-change-quasi-compact}{source}{}{sourcepointer-1}{}}
\newlabel{lemma-base-change-quasi-compact@cref}{{[originallemma][0][]lemma-base-change-quasi-compact}{[1][0][]source}{}{}{}}
\makeatother
\crefname{originallemma}{original lemma}{original lemmas}
\Crefname{originallemma}{Original lemma}{Original lemmas}
\makeatletter
\newlabel{lemma-monomorphism-cartesian-square-inertia}{{lemma-monomorphism-cartesian-square-inertia}{source}{}{sourcepointer-2}{}}
\newlabel{lemma-monomorphism-cartesian-square-inertia@cref}{{[originallemma][0][]lemma-monomorphism-cartesian-square-inertia}{[1][0][]source}{}{}{}}
\makeatother
\crefname{originalproposition}{original proposition}{original propositions}
\Crefname{originalproposition}{Original proposition}{Original propositions}
\makeatletter
\newlabel{proposition-open-stratum}{{proposition-open-stratum}{source}{}{sourcepointer-3}{}}
\newlabel{proposition-open-stratum@cref}{{[originalproposition][0][]proposition-open-stratum}{[1][0][]source}{}{}{}}
\makeatother
\crefname{originalreference}{original reference}{original references}
\Crefname{originalreference}{Original reference}{Original references}
\makeatletter
\newlabel{stacks-properties-lemma-reduced-closed-substack}{{stacks-properties-lemma-reduced-closed-substack}{source}{}{sourcepointer-4}{}}
\newlabel{stacks-properties-lemma-reduced-closed-substack@cref}{{[originalreference][0][]stacks-properties-lemma-reduced-closed-substack}{[1][0][]source}{}{}{}}
\makeatother
\crefname{originalreference}{original reference}{original references}
\Crefname{originalreference}{Original reference}{Original references}
\makeatletter
\newlabel{stacks-properties-lemma-space-locally-quasi-compact}{{stacks-properties-lemma-space-locally-quasi-compact}{source}{}{sourcepointer-5}{}}
\newlabel{stacks-properties-lemma-space-locally-quasi-compact@cref}{{[originalreference][0][]stacks-properties-lemma-space-locally-quasi-compact}{[1][0][]source}{}{}{}}
\makeatother
\title{Every point of a stack lies in a locally closed stratum over which it is a gerbe}
\author{The Stacks Project Authors}
\date{stacks-project:stacks-morphisms.tex:local-snapshot-20261012; GNU FDL 1.2 or later}
\begin{document}
\maketitle
\noindent\textit{Editorial note.} All delivered mathematical text is the author's own, copied from the pinned Stacks Project chapter source sources/stacks-algebraic-stacks-morphisms/pinned-original-source.tex; no statement or proof step was rewritten, reordered or removed.\par\medskip
\noindent\textit{Editorial note.} Preamble provenance: the Stacks Project keeps one single shared preamble for every chapter; this card takes it directly from the staged sources/stacks-algebraic-stacks-morphisms/preamble.tex as three verbatim spans (lines 8-12, 14-18 and 154-217), so no derived file was written and nothing was added or redefined. The wrapper class and the base packages are supplied by the assembler.\par\medskip
\noindent\textit{Editorial note.} Bibliography provenance: neither the counted statement nor its proof contains a citation (this lane rescans the whole span before writing the card), so the card carries no bibliography and no entry was reconstructed.\par\medskip
\noindent\textit{Editorial note.} Selection provenance: this card was opened under the file-based control contract (CONTROL.md), with the mechanical pre-card gate run on its statement and proof interval before the card was written.\par\medskip
\noindent\textit{Allowed standard prerequisites (editorial).}\begin{itemize}
\item The existence of a maximal open stratum on which the inertia is flat and locally of finite presentation. Invoked inside the counted proof (source lines 6849-6860). Produces the open stratum through which the point is a gerbe.
\item The construction of the reduced induced closed substack on the complement. Used inside the counted proof (source lines 6861-6875) as a cross-chapter pointer. Allows the induction step to continue on the closed complement.
\item The base change description of the inertia of a closed substack. Invoked inside the counted proof (source lines 6876-6885). Shows that the complement again satisfies the standing hypotheses.
\end{itemize}\medskip
\ifdefined\hypertarget\hypertarget{sourcepointer-1}{}\fi
\noindent\textit{Source reference.} Original lemma lemma-base-change-quasi-compact: sources/stacks-algebraic-stacks-morphisms/pinned-original-source.tex, source lines 1984--1984. Navigation pointer only: the counted text refers to the Stacks item lemma-base-change-quasi-compact, which is an existing labeled result of this same chapter and is not part of this card. The pointer carries the target's own label as its number and its definition line as its span; its content is not reproduced here.\par\medskip
\ifdefined\hypertarget\hypertarget{sourcepointer-2}{}\fi
\noindent\textit{Source reference.} Original lemma lemma-monomorphism-cartesian-square-inertia: sources/stacks-algebraic-stacks-morphisms/pinned-original-source.tex, source lines 1480--1480. Navigation pointer only: the counted text refers to the Stacks item lemma-monomorphism-cartesian-square-inertia, which is an existing labeled result of this same chapter and is not part of this card. The pointer carries the target's own label as its number and its definition line as its span; its content is not reproduced here.\par\medskip
\ifdefined\hypertarget\hypertarget{sourcepointer-3}{}\fi
\noindent\textit{Source reference.} Original proposition proposition-open-stratum: sources/stacks-algebraic-stacks-morphisms/pinned-original-source.tex, source lines 6746--6746. Navigation pointer only: the counted text refers to the Stacks item proposition-open-stratum, which is an existing labeled result of this same chapter and is not part of this card. The pointer carries the target's own label as its number and its definition line as its span; its content is not reproduced here.\par\medskip
\ifdefined\hypertarget\hypertarget{sourcepointer-4}{}\fi
\noindent\textit{Source reference.} Original reference stacks-properties-lemma-reduced-closed-substack: sources/stacks-algebraic-stacks-morphisms/pinned-original-source.tex, source lines 6854--6854. Navigation pointer only: the counted text refers to the Stacks item stacks-properties-lemma-reduced-closed-substack, which belongs to another chapter of the project and is not part of this card. The pointer carries the target's own label as its number because this chapter's staged snapshot does not contain the other chapter's tag table; its content is not reproduced here.\par\medskip
\ifdefined\hypertarget\hypertarget{sourcepointer-5}{}\fi
\noindent\textit{Source reference.} Original reference stacks-properties-lemma-space-locally-quasi-compact: sources/stacks-algebraic-stacks-morphisms/pinned-original-source.tex, source lines 6871--6871. Navigation pointer only: the counted text refers to the Stacks item stacks-properties-lemma-space-locally-quasi-compact, which belongs to another chapter of the project and is not part of this card. The pointer carries the target's own label as its number because this chapter's staged snapshot does not contain the other chapter's tag table; its content is not reproduced here.\par\medskip
\begin{lemma}
\label{lemma-every-point-in-a-stratum}
Let $\mathcal{X}$ be an algebraic stack such that
$\mathcal{I}_\mathcal{X} \to \mathcal{X}$ is quasi-compact.
Then there exists a well-ordered index set $I$ and for every $i \in I$
a reduced locally closed substack $\mathcal{U}_i \subset \mathcal{X}$ such that
\begin{enumerate}
\item each $\mathcal{U}_i$ is a gerbe,
\item we have $|\mathcal{X}| = \bigcup_{i \in I} |\mathcal{U}_i|$,
\item $T_i = |\mathcal{X}| \setminus \bigcup_{i' < i} |\mathcal{U}_{i'}|$
is closed in $|\mathcal{X}|$ for all $i \in I$, and
\item $|\mathcal{U}_i|$ is open in $T_i$.
\end{enumerate}
We can moreover arrange it so that either (a) $|\mathcal{U}_i| \subset T_i$
is dense, or (b) $\mathcal{U}_i$ is quasi-compact. In case (a), if
we choose $\mathcal{U}_i$ as large as possible (see proof for details), then
the stratification is canonical.
\end{lemma}

\begin{proof}
Let $T \subset |\mathcal{X}|$ be a nonempty closed subset. We are going
to find (resp.\ choose) for every such $T$ a reduced locally closed substack
$\mathcal{U}(T) \subset \mathcal{X}$ with $|\mathcal{U}(T)| \subset T$
open dense (resp.\ nonempty quasi-compact). Namely, by
Properties of Stacks, Lemma
\ref{stacks-properties-lemma-reduced-closed-substack}
there exists a unique reduced closed substack
$\mathcal{X}' \subset \mathcal{X}$ such that $T = |\mathcal{X}'|$.
Note that $\mathcal{I}_{\mathcal{X}'} =
\mathcal{I}_\mathcal{X} \times_\mathcal{X} \mathcal{X}'$ by
Lemma \ref{lemma-monomorphism-cartesian-square-inertia}.
Hence $\mathcal{I}_{\mathcal{X}'} \to \mathcal{X}'$ is
quasi-compact as a base change, see
Lemma \ref{lemma-base-change-quasi-compact}.
Therefore
Proposition \ref{proposition-open-stratum}
implies there exists a dense maximal (see proof proposition)
open substack $\mathcal{U} \subset \mathcal{X}'$
which is a gerbe. In case (a) we set $\mathcal{U}(T) = \mathcal{U}$
(this is canonical) and in case (b) we simply choose a nonempty quasi-compact
open $\mathcal{U}(T) \subset \mathcal{U}$, see
Properties of Stacks, Lemma
\ref{stacks-properties-lemma-space-locally-quasi-compact}
(we can do this for all $T$
simultaneously by the axiom of choice).

\medskip\noindent
Using transfinite recursion we construct for every ordinal $\alpha$ a
closed subset $T_\alpha \subset |\mathcal{X}|$. For $\alpha = 0$
we set $T_0 = |\mathcal{X}|$. Given $T_\alpha$ set
$$
T_{\alpha + 1} = T_\alpha \setminus |\mathcal{U}(T_\alpha)|.
$$
If $\beta$ is a limit ordinal we set
$$
T_\beta = \bigcap\nolimits_{\alpha < \beta} T_\alpha.
$$
We claim that $T_\alpha = \emptyset$ for all $\alpha$
large enough. Namely, assume that $T_\alpha \not = \emptyset$
for all $\alpha$. Then we obtain an injective map from the class
of ordinals into the set of subsets of $|\mathcal{X}|$ which is a
contradiction.

\medskip\noindent
The claim implies the lemma. Namely, let
$$
I = \{\alpha \mid \mathcal{U}_\alpha \not = \emptyset \}.
$$
This is a well-ordered set by the claim. For $i = \alpha \in I$ we set
$\mathcal{U}_i = \mathcal{U}_\alpha$. So $\mathcal{U}_i$ is a reduced
locally closed substack and a gerbe, i.e., (1) holds. By construction
$T_i = T\alpha$ if $i = \alpha \in I$, hence (3) holds. Also, (4) and
(a) or (b) hold by our choice of $\mathcal{U}(T)$ as well. Finally, to see
(2) let $x \in |\mathcal{X}|$. There exists a smallest ordinal $\beta$
with $x \not \in T_\beta$ (because the ordinals are well-ordered).
In this case $\beta$ has to be a successor ordinal by the definition
of $T_\beta$ for limit ordinals. Hence $\beta = \alpha + 1$ and
$x \in |\mathcal{U}(T_\alpha)|$ and we win.
\end{proof}

\end{document}
